% Fragmento integrable; no documento autónomo y no orden de compilación.
% Fecha: 2026-09-30. No modifica fuentes publicadas.
% Procedencia: reunión de RETROACCION_PCH_Y_RESTRICCIONES_VARIACIONALES.md,
% PRECUANTIZACION_PCH_Y_FIDELIDAD.md y CIERRE_CANONICO_TOTAL_Y_REPRESENTACION.md.
% Prerrequisitos de inserción: núcleo común APP--TRIT--TPK; VIII/30d,
% VIII/31 y VIII/33; lectores de acción, gamma y acoplamientos.
% Los nuevos labels usan exclusivamente el prefijo hmtcuatro:.
\section{Action conjointe, rétroaction et clôture canonique des contraintes}
\label{hmtcuatro:accion-retroaccion-restricciones}

La séquence causale de ce développement est
\[
 \mathrm{APP}\longrightarrow\mathrm{TRIT}\longrightarrow\mathrm{TPK}
 \longrightarrow\text{état enrichi}
 \longrightarrow\text{structure discrète conjointe du continu}.
\]
Les feuillets, l’orientation, les résidus, la retenue, les routes, la frontière et la mémoire de cette construction sont conservés. L’échelle d’action et les couplages sont des sorties HMT reçues avant l’emploi du langage variationnel et symplectique. Ils ne sont pas sélectionnés par des valeurs cibles ni par la conclusion de courbure à démontrer. Le retour de phase ne réinitialise pas non plus l’état enrichi.

La réalisation géométrique précédemment construite fournit une cotétrade non dégénérée \(e\), sa connexion métrique et le relèvement spinoriel du transport. Ici, cette géométrie varie avec les champs matériels. Nous réunirons en une seule démonstration l’action commune, sa réduction torsionnelle, la transformation de Legendre, le caractère de première classe, la propagation des contraintes et la platitude caractéristique de la connexion canonique issue de la même action. Cette dernière réalisation n’est pas identifiée par son nom à un autre hamiltonien en interaction.

\subsection{Structure de départ et action à source complète}
\label{hmtcuatro:datos-accion-total}

Soit \(M\) une région lisse orientée de dimension quatre, munie d’une structure de spin, de signature \((-+++)\) et d’un feuilletage spatial régulier. La cotétrade prend ses valeurs dans la droite de longueur de la réalisation reçue. Nous écrivons
\[
 \Sigma^{IJ}=e^I\wedge e^J,\qquad
 F_\omega=\mathrm d\omega+\omega\wedge\omega,\qquad
 \mathsf P_\gamma=\star+\gamma^{-1}I,\qquad \star^2=-I.
\]
L’étoile de cette formule agit sur les bivecteurs internes ; ce n’est pas la dualité extérieure des formes. Pendant la variation, les sections reçues \(c>0,\hbar>0,G>0,\gamma\in\mathbb R\setminus\{0\}\) et \(\Lambda\) restent fixes. Avec \(\kappa=8\pi G/c^4\), l’action totale est
\begin{equation}
 \begin{split}
 S_{\rm tot}={}&\frac1{2\kappa c}
   \int_M\Sigma_{IJ}\wedge(\mathsf P_\gamma F_\omega)^{IJ}
 -\frac{\Lambda}{\kappa c}\int_M\operatorname{vol}_e\\
 &+S_{\rm YM}[e,A]+S_H[e,A,H]
   +S_W[e,\omega,A,\Psi]+S_Y[e,H,\Psi]+S_{\partial}.
 \end{split}
 \label{hmtcuatro:accion-total}
\end{equation}
Le symbole \(H\) sans indice désigne ici le champ de Higgs, non un hamiltonien. Toutes les contractions spatiotemporelles de la matière utilisent la même \(e\). Le problème de bord est celui reçu de la variation de Cartan--Holst : variations intérieures, données de bord fixes compatibles ou complétion annulant le même terme de frontière.

Pour fixer le contenu de \eqref{hmtcuatro:accion-total}, les connexions internes agissent sur le fibré chiral
\[
 (S_L\otimes E_L)\oplus(S_R\otimes E_R).
\]
Ce fibré n’est pas remplacé par un couplage vectoriel. Dans la normalisation entière \(y=6Y\), les multiplets reçus ont pour poids \(1,4,-2,-3,-6\), respectivement pour \(q_L,u_R,d_R,\ell_L,e_R\) ; le doublet \(H\) a pour poids \(3\). La couleur agit sur le facteur \(\mathbb C^3\) des quarks et trivialement sur les leptons. Le facteur de familles et ses applications \(Y_u,Y_d,Y_e\) sont conservés. Ce bloc n’ajoute pas de partenaire droit à un multiplet qui n’en possède pas.

Avec \(\widetilde H=i\sigma_2\overline H\), l’application matérielle est
\begin{equation}
 \begin{split}
 \mathsf Y_q(H)(u,d)&=
       \widetilde H\otimes Y_u u+H\otimes Y_d d,\\
 \mathsf Y_\ell(H)e&=H\otimes Y_e e .
 \end{split}
 \label{hmtcuatro:yukawa}
\end{equation}
L’identité sur le facteur de couleur est sous-entendue à la première ligne. Le terme d’action est \(-(\bar\Psi_L\mathsf Y(H)\Psi_R+\mathrm{h.c.})\), avec les sections dimensionnelles reçues. L’équivariance n’est pas seulement nominale : pour \(U\in SU(2)\), \(i\sigma_2\overline U=U i\sigma_2\), de sorte que \(\widetilde H\mapsto U\widetilde H\). Les poids des trois colonnes sont \(-3+4=1\), \(3-2=1\) et \(3-6=-3\). La couleur est entrelacée par l’identité et les applications de famille commutent avec l’action interne. Par conséquent, pour chaque transformation interne \(u\),
\begin{equation}
 \mathsf Y(uH)\rho_R(u)=\rho_L(u)\mathsf Y(H).
 \label{hmtcuatro:equivariancia-yukawa}
\end{equation}
Si la réalisation matérielle précédente contient en outre d’autres multiplets, leurs applications sont incorporées avec leur identité d’équivariance propre ; cette preuve n’invente pas de nouveaux états.

Dans une carte d’unités \(\hbar=c=1\), la densité matérielle qui résume ces actions a la forme
\begin{equation}
 \begin{split}
 \mathcal L_m={}&-\frac14 F_{B,\mu\nu}F_B^{\mu\nu}
 -\frac14\sum_a F_{W,\mu\nu}^aF_W^{a,\mu\nu}
 -\frac14\sum_a F_{C,\mu\nu}^aF_C^{a,\mu\nu}\\
 &-(\nabla_\mu H)^\dagger\nabla^\mu H-V(H)
 +\mathcal L_W(\Psi_L,\nabla_L)+\mathcal L_W(\Psi_R,\nabla_R)\\
 &-\bigl(\bar\Psi_L\mathsf Y(H)\Psi_R+\mathrm{h.c.}\bigr).
 \end{split}
 \label{hmtcuatro:densidad-material}
\end{equation}
Les courbures incluent leurs couplages ; ceux-ci et les normalisations des composantes sont ceux de leurs documents sources. L’écriture dans cette carte n’élimine pas la section d’action et ne recalibre pas les coefficients. En particulier, aucune nouvelle valeur numérique du couplage fort n’est introduite ici. La connexion de spin de \(\nabla_L,\nabla_R\) est celle de \(e,\omega\) ; les connexions \(B,W,C\) agissent dans les représentations qui viennent d’être décrites. La constante de \(V\) est conservée, car sa variation de volume appartient à la source gravitationnelle.

La partie spinorielle conserve les conventions explicites du développement précédent : \(\{g^I,g^J\}=2\eta^{IJ}I\), \(\mathbb B_{IJ}=\frac14[g_I,g_J]\), \(A_D=-ig^0\), \(\bar\psi=\psi^\dagger A_D\) et \(D\psi=\mathrm d\psi+\frac12\omega^{IJ}\mathbb B_{IJ}\psi+A\psi\). Sa cinétique symétrique est
\[
 \frac{\hbar}{2}
 [\bar\psi g^I D\psi-(D\bar\psi)g^I\psi]\wedge\eta_I,
 \qquad \eta_I=\iota_{E_I}\operatorname{vol}_e .
\]
La somme porte sur les champs réellement présents, avec leurs projections chirales, non sur une duplication vectorielle du catalogue. Les termes de Yukawa sont algébriques et ne dépendent pas de \(\omega\). Le courant total est fixé par
\begin{equation}
 \delta_\omega S_m=-\frac12\int_M\delta\omega^{IJ}\wedge\sigma_{IJ},
 \qquad
 \sigma_{JK}=\sum_s-\frac{\hbar}{2}
   \bar\psi_s\{g^I,\mathbb B_{JK}\}\psi_s\,\eta_I.
 \label{hmtcuatro:corriente-total}
\end{equation}

\subsection{Une démonstration de compatibilité, de propagation et de clôture}
\label{hmtcuatro:demostracion-unica}

\begin{theorem}[Compatibilité canonique de l’action conjointe]
\label{hmtcuatro:teorema-canonico-total} Considérons \eqref{hmtcuatro:accion-total} sur le domaine non dégénéré précédent, avec une matière affine en la contorsion, et dans une carte régulière de la transformation de Legendre restreinte. Conservons les contraintes fermioniques du premier ordre et éliminons celles de seconde classe au moyen de leur crochet de Dirac. Pour des paramètres à support intérieur, ou avec le bord compatible sans charges impropres, les propriétés suivantes sont conjointement satisfaites : l’équation métrique possède la source matérielle totale de cette même action ; ses générateurs de difféomorphismes et de jauge sont de première classe ; les contraintes se propagent sur toute solution couplée régulière ; et le potentiel canonique total induit une connexion préquantique plate dans les directions caractéristiques de la surface régulière des contraintes.

L’assertion porte sur cette action et cette réalisation canonique. Elle n’identifie pas automatiquement sa connexion à un hamiltonien en interaction distinct et ne prouve pas sa limite spatiale--chirale.
\end{theorem}

\begin{proof}
\paragraph{1. Élimination torsionnelle sans duplication.}
Écrivons \(\omega=\mathring\omega(e)+\mathfrak k\). Le développement \(F_\omega=F_{\mathring\omega}
+D_{\mathring\omega}\mathfrak k+\mathfrak k\wedge\mathfrak k\) et \(D_{\mathring\omega}\Sigma=0\) séparent la contribution de bord
\[
 S_{\partial,\mathfrak k}
 =\frac1{2\kappa c}\int_{\partial M}
                  \Sigma_{IJ}\wedge(\mathsf P_\gamma\mathfrak k)^{IJ}
\]
de la densité algébrique
\begin{equation}
 \mathcal Q_{e,\gamma}(\mathfrak k)
       -\frac12\mathfrak k^{IJ}\wedge\sigma_{IJ},\qquad
 \mathcal Q_{e,\gamma}(\mathfrak k)=\frac1{2\kappa c}
       \Sigma_{IJ}\wedge
       \mathsf P_\gamma(\mathfrak k\wedge\mathfrak k)^{IJ}.
 \label{hmtcuatro:cuadratica-contorsion}
\end{equation}
La variation originale donne \(D_\omega(\mathsf P_\gamma\Sigma)=\kappa c\,\sigma\). L’inverse de Holst et l’inverse torsion--contorsion du développement précédent sont inversibles pour le corepère et \(\gamma\) déclarés. Ils déterminent ainsi une unique \(\mathfrak k_*[e,\sigma]\), linéaire en \(\sigma\). La variation radiale de \eqref{hmtcuatro:cuadratica-contorsion} en son point stationnaire donne \(2\mathcal Q(\mathfrak k_*)=\frac12\mathfrak k_*^{IJ}\wedge\sigma_{IJ}\). Par conséquent,
\begin{equation}
 S_{\rm red}=\frac1{2\kappa c}\int_M(R-2\Lambda)\operatorname{vol}_e
       +S_m^{\rm LC}
       -\frac14\int_M\mathfrak k_*^{IJ}\wedge\sigma_{IJ}
       +S_{\partial,\rm red}.
 \label{hmtcuatro:accion-reducida}
\end{equation}
La première identité de Bianchi annule le terme de Holst de Levi--Civita dans cette écriture. La règle de la chaîne conserve les équations de \(e,A,H,\Psi\), car le coefficient de \(\delta\mathfrak k_*\) s’annule par stationnarité. Comme \(\sigma=\sum_s\sigma_s\), la contribution effective inclut \(\mathfrak k_*[\sigma_s]\wedge\sigma_t\) pour \(s\ne t\). Elle n’est pas remplacée par des carrés d’espèces isolées, et la même \(\mathfrak k\) n’est pas conservée simultanément comme variable indépendante.

\paragraph{2. Source complète et rétroaction.}
Définissons le tenseur d’action par la variation de tous les termes matériels de \eqref{hmtcuatro:accion-reducida} :
\begin{equation}
 \delta_g S_{m,\rm red}=-\frac12\int_M
       \mathcal T^S_{\mu\nu}\,\delta g^{\mu\nu}\operatorname{vol}_e,
 \qquad
 E_{\mu\nu}:=G_{\mu\nu}+\Lambda g_{\mu\nu}
                   -\kappa c\,\mathcal T^S_{\mu\nu}.
 \label{hmtcuatro:fuente-completa}
\end{equation}
L’équation métrique est \(E_{\mu\nu}=0\). Le facteur \(\kappa c\) correspond à la cotétrade de longueur ; si \(T^{\rm phys}=c\mathcal T^S\), on écrit \(\kappa T^{\rm phys}\). On fait également varier les champs matériels et les connexions internes. Les champs déterminent donc la source, et la géométrie qui résout l’équation détermine à son tour leurs connexions, leur volume et leur évolution. Il s’agit d’une rétroaction, non d’une évolution sur un fond fixé. La variation inclut le volume du potentiel de Higgs et la dépendance métrique du courant torsionnel, ainsi que ses termes croisés.

\paragraph{3. Transformation de Legendre restreinte et réduction canonique.}
Le potentiel canonique total s’obtient à partir de \(\delta L=\mathcal E\delta z+\mathrm d\theta(\delta z)\), en intégrant \(\theta\) sur la feuille spatiale et en effectuant la réduction de seconde classe. Sa partie géométrique avant cette réduction est
\[
 \theta_{\rm geom}(\delta)=
   \frac1{2\kappa c}(\mathsf P_\gamma\Sigma)_{IJ}
                       \wedge\delta\omega^{IJ}.
\]
La décomposition \(F_{0a}=\dot\omega_a-D_a\omega_0\), avec la même décomposition des actions matérielles, produit
\begin{equation}
 S_{\rm can}=\int d\tau\,
 \bigl[\Theta_{\rm red}(\dot z)-\mathcal C[N]-\mathcal D[\xi]
       -\mathcal G_{\rm L}[\lambda]-\mathcal G_{\rm int}[\chi]\bigr]
       +S_{\partial,\rm red}.
 \label{hmtcuatro:legendre-total}
\end{equation}
Les vitesses absentes du lapse, du shift ou des connexions temporelles ne sont pas artificiellement inversées.

La contorsion algébrique possède les contraintes \(\pi_{\mathfrak k}=0\) et \(\partial\mathcal Q/\partial\mathfrak k-\sigma/2=0\). Leur matrice de crochets possède le bloc
\[
 \begin{pmatrix}0&-M\\ M^T&B\end{pmatrix},
 \qquad M=\frac{\partial^2\mathcal Q}{\partial\mathfrak k^2}.
\]
L’inverse torsionnel rend \(M\) inversible, et donc cette matrice. Ce sont des contraintes de seconde classe. On utilise
\begin{equation}
 \{F,G\}_D=\{F,G\}
       -\{F,\chi_a\}(C^{-1})^{ab}\{\chi_b,G\}.
 \label{hmtcuatro:corchete-dirac}
\end{equation}
Pour les variables qui ne dépendent pas de la paire contorsion--moment, le bloc inférieur droit de \(C^{-1}\) est nul. Cette élimination ne fait apparaître aucun crochet supplémentaire entre elles. Les contraintes spinorielles du premier ordre sont conservées séparément ; dans l’espace de phase gradué, on emploie le crochet gradué et les générateurs pairs. \(\Theta_{\rm red}\) désigne le résultat complet, non sa seule partie géométrique.

L’inversion de Legendre dans les directions régulières restitue la même action. Par exemple, \(N[p^2/(2w)+wV+H_{\rm grad}]\), avec \(w=\sqrt{\det q}\), donne \(p=(w/N)D_\tau\phi\), et par substitution
\[
 L_H=\frac{w}{2N}|D_\tau\phi|^2-NwV-NH_{\rm grad}.
\]
\(D_\tau\) conserve le shift et la connexion temporelle ; \(w\) n’est pas figé. Pour \(K_{ij}=(\dot q_{ij}-\mathcal L_\xi q_{ij})/(2N)\), la contribution gravitationnelle au moment et son inverse sont
\begin{equation}
 \Pi^{ij}=\frac{\sqrt q}{2\kappa c}(K^{ij}-q^{ij}K),\qquad
 K_{ij}=\frac{2\kappa c}{\sqrt q}
                  (\Pi_{ij}-\tfrac12q_{ij}\Pi).
 \label{hmtcuatro:legendre-geometrica}
\end{equation}
Il en résulte
\[
 \mathcal C_{\rm grav}
 =\frac{2\kappa c}{\sqrt q}
       (\Pi^{ij}\Pi_{ij}-\tfrac12\Pi^2)
  -\frac{\sqrt q}{2\kappa c}(R^{(3)}-2\Lambda).
\]
Si la carte des spineurs apporte un décalage matériel du moment total \(p\), on utilise \(\Pi=p-p_m\) ; \(p_m\) n’est pas écarté. Le caractère indéfini de la forme de DeWitt n’empêche pas cet inverse algébrique. On n’exige pas d’inverse des vitesses pour la cinétique fermionique du premier ordre.

\paragraph{4. Identité de Noether--Legendre et caractère de première classe.}
La coïncidence des générateurs avec les contraintes de la même action peut être constatée avant d’imposer ses équations. Posons \(a=(2\kappa c)^{-1}\), \(B=\mathsf P_\gamma\Sigma\), \(u^I=\iota_\zeta e^I\) et \(\lambda^{IJ}=\iota_\zeta\omega^{IJ}\). De \(\mathcal L_\zeta\omega=\iota_\zeta F+D_\omega\lambda\), on obtient
\begin{equation}
 \begin{split}
 j_\zeta^{\rm geom}
 ={}&\mathrm d(a B_{IJ}\lambda^{IJ})
   -a(D_\omega B_{IJ})\lambda^{IJ}\\
 &-2a u^Ie^J\wedge(\mathsf P_\gamma F)_{IJ}
       +2a\Lambda u^I\eta_I .
 \end{split}
 \label{hmtcuatro:noether-legendre}
\end{equation}
La partie matérielle ajoute exactement ses contributions aux équations de \(e,\omega,A\), ses générateurs verticaux et son bord. Décomposer la même expression sur une feuille spatiale donne \eqref{hmtcuatro:legendre-total}. La réduction stationnaire et la réduction de Dirac transportent ces transformations au lieu de les remplacer par les moments d’une action extérieure.

Avec la convention canonique usuelle pour les contraintes, les crochets totaux sont, modulo les générateurs de Gauss conservés,
\begin{equation}
 \begin{aligned}
 \{\mathcal D[\xi],\mathcal D[\eta]\}_D
    &\simeq\mathcal D[[\xi,\eta]],\\
 \{\mathcal D[\xi],\mathcal C[N]\}_D
    &\simeq\mathcal C[\mathcal L_\xi N],\\
 \{\mathcal C[N],\mathcal C[M]\}_D
    &\simeq\mathcal D[
       q^{ij}(N\partial_jM-M\partial_jN)\partial_i].
 \end{aligned}
 \label{hmtcuatro:algebra-restricciones}
\end{equation}
En effet, les transformations tangentielles ont le crochet de Lie et transportent le lapse comme un scalaire. Pour les transformations normales, différencier leur orthogonalité aux vecteurs tangents et \(g(n,n)=-1\) détermine la variation tangentielle de \(n\) ; avec \(K_{ij}\) de \eqref{hmtcuatro:legendre-geometrica}, on obtient \(\delta_N n=\operatorname{grad}_qN\). Le commutateur des déformations, converti en crochet canonique avec ces conventions, donne la dernière ligne de \eqref{hmtcuatro:algebra-restricciones}. L’usage du transport covariantisé conserve en outre les rotations de repère et de jauge. Les transformations sont tangentes à la surface des contraintes de l’action invariante, d’où le caractère de première classe. Les fonctions \(q^{ij}\) dépendent de l’état et ne sont pas figées. Avec un support intérieur, aucune charge centrale de bord n’est ajoutée ; un générateur impropre de charge non nulle n’est pas déclaré contrainte nulle.

\paragraph{5. Propagation dans la géométrie qui répond à la matière.}
L’invariance par difféomorphismes de \(S_{m,\rm red}\), évaluée sur ses équations matérielles, donne par intégration par parties \(\mathring\nabla^\mu\mathcal T^S_{\mu\nu}=0\). L’identité utilise le tenseur complet de \eqref{hmtcuatro:fuente-completa} ; elle ne postule pas la conservation séparée de chaque contribution. Bianchi implique alors \(\mathring\nabla_\mu E^{\mu\nu}=0\), avant même d’imposer toutes les équations métriques.

Dans une carte gaussienne locale, \(ds^2=-d\tau^2+q_{ij}dx^idx^j\), imposons les équations d’évolution \(E_{ij}=0\) et posons \(C=E_{00}\), \(C_i=E_{0i}\), \(K_{ij}=\dot q_{ij}/2\), \(K=q^{ij}K_{ij}\). Ici, \(K\) est seulement la trace de la seconde forme fondamentale. Les deux composantes de l’identité de divergence sont
\begin{equation}
 \partial_\tau C-D_iC^i+KC=0,\qquad
 \partial_\tau C_i+KC_i=0.
 \label{hmtcuatro:propagacion}
\end{equation}
Pour vérifier les signes, on utilise \(E^{00}=C\), \(E^{0i}=-C^i\), \(E^{ij}=0\), \(\Gamma^i{}_{0j}=K^i{}_j\) et \(\Gamma^\mu{}_{\mu0}=K\). En abaissant l’indice de l’équation spatiale, \(\dot q_{ij}=2K_{ij}\) annule les deux termes de connexion. Si initialement \(C_i=0\), son équation homogène maintient \(C_i=0\) ; la première devient \(\dot C+KC=0\) et conserve \(C=0\). \(q\) et \(K\) sont ceux de la solution couplée. Le choix gaussien prouve l’assertion locale ; sa forme tensorielle permet de changer le lapse et le shift. De même, l’identité de jauge de Noether relie la divergence covariante de l’équation de jauge aux équations matérielles et propage Gauss lorsque les équations spatiales sont satisfaites. On n’en infère ni existence globale ni absence de singularités pour toute solution.

\paragraph{6. La même action induit la clôture caractéristique.}
Soit \(\mathcal P_{\rm red}\) la carte canonique régulière obtenue et \(\mathcal Z=\{C_A=0\}\) sa surface régulière de contraintes totales, Gauss compris. Nous utilisons pour cette surface un symbole distinct de \(\Sigma^{IJ}\). Le potentiel \(\Theta=\Theta_{\rm red}\) de \eqref{hmtcuatro:legendre-total} est conservé, termes matériels et mixtes compris. Nous fixons
\begin{equation}
 \Omega=-\mathrm d_{\mathcal P}\Theta,\qquad
 \iota_{X_f}\Omega=\mathrm d_{\mathcal P}f,\qquad
 \{f,g\}_D=\Omega(X_f,X_g).
 \label{hmtcuatro:convenciones-simpecticas}
\end{equation}
Dans une paire canonique, \(\Theta=p\,\mathrm dq\) donne \(X_f=f_p\partial_q-f_q\partial_p\), \(\{q,p\}_D=1\), \(X_f(g)=-\{f,g\}_D\) et \([X_f,X_g]=-X_{\{f,g\}_D}\). Dans l’espace de phase gradué, cette écriture s’applique aux fonctions et contraintes paires pertinentes.

Sur la droite de phase trivialisée au-dessus de cette carte, nous construisons
\begin{equation}
 \nabla_X=X-\frac{i}{\hbar}\Theta(X)
       =\delta_X+\frac{i}{\hbar}H_X^{\rm can},
 \qquad H_X^{\rm can}=-\Theta(X).
 \label{hmtcuatro:conexion-canonica}
\end{equation}
On ne choisit pas de conjugaison arbitraire pour la rendre plate : son coefficient provient du potentiel de l’action. Pour tous champs \(X,Y\) de la carte, le calcul donne
\begin{equation}
 \begin{aligned}
 \mathcal F^{\rm can}_{X,Y}
 &=XH_Y^{\rm can}-YH_X^{\rm can}-H_{[X,Y]}^{\rm can}
       +\frac{i}{\hbar}[H_X^{\rm can},H_Y^{\rm can}]\\
 &=-X\Theta(Y)+Y\Theta(X)+\Theta([X,Y])\\
 &=-\mathrm d_{\mathcal P}\Theta(X,Y)=\Omega(X,Y).
 \end{aligned}
 \label{hmtcuatro:curvatura-calculada}
\end{equation}
Les \(H_X^{\rm can}\) sont des multiplicateurs de la droite, de sorte que leur commutateur est nul ; les dérivées des coefficients et des champs ne sont pas supprimées.

Le caractère de première classe déjà prouvé signifie \(\{C_A,C_B\}_D=f_{AB}{}^D C_D\). Pour \(V\) tangent à \(\mathcal Z\), \(\Omega(X_{C_A},V)=\mathrm d_{\mathcal P}C_A(V)=0\). Les champs des contraintes sont donc caractéristiques. Sous la régularité déclarée, ils engendrent la distribution caractéristique de la surface coïsotrope. La tensorialité de \(\Omega\) étend le calcul à toutes leurs combinaisons lisses et prouve
\begin{equation}
 \left.\mathcal F^{\rm can}_{X,Y}\right|_{\mathcal Z}=0
 \quad\hbox{pour }X,Y\hbox{ caractéristiques}.
 \label{hmtcuatro:cierre-caracteristico}
\end{equation}
Pour deux lapses, \([X,Y]\) inclut la déformation tangentielle et les actions de Gauss de \eqref{hmtcuatro:algebra-restricciones}. Si l’on passe au quotient de jauge, l’égalité s’exprime modulo ces actions verticales. On n’affirme pas que \(\Omega\) soit nul sur tout l’espace de phase ni que toute holonomie globale d’une feuille soit triviale.

\paragraph{7. Expression préquantique de la même clôture.}
La réalisation différentielle associée à \eqref{hmtcuatro:conexion-canonica}, sur des sections lisses où les produits sont définis, est
\begin{equation}
 Q_{\rm pre}(f)=-i\hbar\nabla_{X_f}+f
            =-i\hbar X_f+f-\Theta(X_f).
 \label{hmtcuatro:operador-precuantico}
\end{equation}
Avec \(R^\nabla=(i/\hbar)\Omega\), le développement du commutateur et \(X_f(g)=-\{f,g\}_D\) donnent
\begin{equation}
 [Q_{\rm pre}(f),Q_{\rm pre}(g)]
           =i\hbar Q_{\rm pre}(\{f,g\}_D).
 \label{hmtcuatro:representacion-precuantica}
\end{equation}
La règle du produit qui conserve les fonctions de structure est
\begin{equation}
 Q_{\rm pre}(aC)
   =aQ_{\rm pre}(C)+C\bigl(Q_{\rm pre}(a)-a\bigr).
 \label{hmtcuatro:funciones-estructura}
\end{equation}
Elle découle de \(X_{aC}=aX_C+CX_a\) et de \eqref{hmtcuatro:operador-precuantico}. Le second terme n’est pas éliminé hors de \(\mathcal Z\). Sur \(\mathcal Z\), \(Q_{\rm pre}(C_A)=-i\hbar\nabla_{X_{C_A}}\) ; ainsi, \eqref{hmtcuatro:representacion-precuantica} et \eqref{hmtcuatro:cierre-caracteristico} sont les expressions opératorielle et géométrique de la même clôture canonique, non deux incorporations indépendantes de la gravité.
\end{proof}

\subsection{Domaine du résultat et comparaison des réalisations}
\label{hmtcuatro:alcance-canonico}

La preuve réunit une action à gravité dynamique et source conjointe, sa transformation de Legendre et sa connexion caractéristique. Le qualificatif \emph{arbitraire} pour les déformations se rapporte aux lapses et déplacements admissibles qui engendrent des directions caractéristiques dans cette région régulière ; non à tout vecteur de phase, tout bord ou toute solution singulière. La platitude locale conserve l’éventuelle monodromie globale et n’identifie pas deux histoires TPK du seul fait qu’elles ont les mêmes extrémités.

Le coefficient \(H_X^{\rm can}\) d’une connexion de droite n’est pas par définition l’opérateur d’horloge et de réponse \(H_{\rm clk}+D^*WD\), ni son extension en interaction. \(Q_{\rm pre}(C)\) inclut en outre une dérivation sur l’espace de phase ; ce n’est pas non plus un autre nom de ces opérateurs. Pour identifier une réalisation concrète au moyen d’une application \(I\), l’égalité typée est requise
\[
 \nabla_N^{\rm op}I=I\nabla_N^{\rm can}
\]
sur leurs domaines effectifs. Si \(I\) dépend de la géométrie, sa dérivée fait partie de cette égalité.

La distinction dispose d’un contrôle exact : dans la carte \(\Theta=p\,\mathrm dq\),
\begin{equation}
 Q_{\rm pre}(p^2/2)=-i\hbar p\,\partial_q-p^2/2.
 \label{hmtcuatro:control-polarizacion}
\end{equation}
En agissant sur la fonction \(1\), il produit \(-p^2/2\). Il ne préserve donc pas par simple restriction la polarisation des fonctions indépendantes de \(p\) et n’y est pas \(-\hbar^2\partial_q^2/2\). Le contrôle ne réfute pas l’hamiltonien reçu : il empêche de remplacer une représentation par une autre sans son application.

Aucune mesure de Liouville de dimension infinie n’a été supposée et l’autoadjonction ne se déduit pas des identités différentielles précédentes. Si les inclusions de raffinement entrelacent les connexions et conservent un cœur commun, leurs courbures s’entrelacent par composition ; une limite exige en outre ses propres contrôles de domaine et de convergence. Cette condition n’est pas présentée ici comme vérifiée pour l’opérateur en interaction et la limite spatiale--chirale conjointe. Les constructions quantiques antérieures sont conservées ; le résultat présent précise exactement la clôture canonique qui vient d’être démontrée.
